% Paul E. West and Sean T. Hayes

\documentclass[14pt, aspectratio=1610, hyperref={pdfpagelayout=SinglePage}]{beamer}
% \documentclass[14pt, aspectratio=1610, hyperref={pdfpagelayout=SinglePage}, handout]{beamer}
\usepackage{csu-theme}

% for aligning items in an itemize (eqmakebox)
\usepackage{eqparbox}

\title{\Huge{}Classes, Composition, \& Inheritance in {\ttfamily\bfseries C++}}
\subtitle{CSCI 315: Data Structure Analysis\\Chapter 10 \& 11}
%\author{}
\institute{Charleston Southern University}

%% Comment out date to use default (today)
% \date{January 28, 2025}
\subject{Software Programming}
%\keywords{}

\begin{document}

\begin{frame}
  \titlepage
\end{frame}

\section{Intro to Classes}


\begin{frame}[fragile]{Syntax of Class}
\begin{itemize}
\item
  The general syntax for defining a class is
\begin{lstlisting}
class ClassIdentifier
{
  // Class members (variables and functions);
}; // <- Don't forget the semicolon.
\end{lstlisting}
\item<2->
  A class member can be a variable or a function.
\end{itemize}
\end{frame}

\begin{frame}[fragile]{Syntax of Class}
  \begin{itemize}
  \item<+->
    If a member of a class is a variable,
    \begin{itemize}
    \item
      It is declared like any other variable.
    \item
      You cannot initialize a variable when you declare it.
    \end{itemize}
  \item<+->
    If a member of a class is a function,
    \begin{itemize}
    \item
      A function prototype declares that member.
    \item
      Function members can \emph{directly} access any member of the class.
    \end{itemize}
  \end{itemize}
\end{frame}

\begin{frame}{Class Organization}
  Access restrictions for members are set with the following keywords.
  \begin{itemize}
  \item \lstinline{public}: accessible outside the class.
  \item \lstinline{private}: cannot be accessed outside the class.
  \item \lstinline{protected}: access for the class and its subclasses.
  \item \lstinline{friend}: grant member-level access to non-member functions or classes. Use sparingly (or not at all); this reduces encapsulation and increase coupling.
  \end{itemize}
\end{frame}


\begin{frame}{Class Organization}
  \begin{itemize}
  \item <+->
    In Java, each member is prefixed with a keyword.
  \item <+->
    In C++,
    \begin{itemize}
    \item Each class can have public, private, and protected sections listing members with that access.
    \item The \lstinline{friend} keyword is placed before a function/class name to indicate its access level.
    \end{itemize}
  \end{itemize}
\end{frame}

\begin{frame}[fragile]{A \textbf{\texttt{struct}} versus a \textbf{\texttt{class}}}
\begin{itemize}[<+->]
\item
  By default, members of a \lstinline{struct} are public.\\
  Use the \lstinline{private} specifier to make a member private.
\item
  By default, the members of a \lstinline{class} are private.
\item
  Classes and structs have the same capabilities.
\item
  Use \lstinline{struct} if all member variables of a class are public and especially if there are no member functions (``plain old data'').
\end{itemize}
\end{frame}

\section{Member Functions}

\begin{frame}{Member Functions}
\begin{itemize}
\item \emph{Methods} are member functions, which define object behavior.
\item Including only prototypes in class declaration keeps the class smaller and hides the implementation.
\end{itemize}
\end{frame}

\begin{frame}[fragile]{Example Class}
\begin{columns}[t]
\column{0.45\textwidth}
\texttt{AClass.hpp}
\begin{lstlisting}[basicstyle=\small\ttfamily, xleftmargin=0pt]
#ifndef A_CLASS
#define A_CLASS
class AClass {
  public:
    // Methods
    int getSize() const;
    void size(int size);

  private:
    int mSize;
};
#endif
\end{lstlisting}
\column{0.5\textwidth}
\texttt{AClass.cpp}
\begin{lstlisting}[basicstyle=\small\ttfamily, xleftmargin=0pt]
#include "AClass.hpp"
int AClass::getSize() const
{
  return mSize;
  // "this" is implied.
  // return this->mSize;
}

void AClass::size(int size)
{
  mSize = size;
}
\end{lstlisting}
\end{columns}
\end{frame}

\begin{frame}[fragile]{Class File Organization}
  \begin{itemize}
  \item
    Typically, each class declaration is in its own \texttt{.hpp} file.
    \begin{itemize}
    \item
      Avoid including this code multiple times in a program,
      by using preprocessor directives to check for a unique definition.\\[0.5em]
      \texttt{MyClass.hpp}
      \begin{lstlisting}[basicstyle=\small\ttfamily, xleftmargin=0pt]
#ifndef MY_CLASS
#define MY_CLASS
//...
#endif\end{lstlisting}
    \end{itemize}
  \item
    The function definitions go in the .cpp file with the matching name and will include the header.\\
      e.g., the top of \texttt{MyClass.cpp} has \lstinline{#include "MyClass.hpp"}.
  \end{itemize}
  \end{frame}

\begin{frame}{Types of Class Methods (Member Functions)}
\begin{description}[Accessor:]
\item[\textbf{\textcolor{white}{Mutator}}:]<+->
  modifies the value(s) of member variable(s).
\item[\textbf{\textcolor{white}{Accessor}}:]<+->
  only accesses the value(s) of member variable(s).
\item[\textbf{\textcolor{white}{Constant}}:]<+->
  cannot modify the class's member variables.
  \begin{itemize}
  \item
    The function heading ends with \lstinline{const}.
  \item
    Most accessors should be constant.
  \end{itemize}
\item[\textbf{\textcolor{white}{Static}}:]<+->
  does \textbf{not} have access to non-static members.\\Does not have access to the \texttt{this} pointer.
\end{description}
\end{frame}



\begin{frame}{Types of Member Functions}
\begin{itemize}
\item<1->
  Instance methods -- accessible through an instantiated object.
\item<2->
  \emph{Static} methods -- accessible using the class name and the scope resolution operator (\lstinline{::}).
  \begin{itemize}
    \item Use the keyword \lstinline{static} to declare a function or variable of a class as static.
    \item No object is needed to access a static member.
    \item Static methods do not have access not no-static members.
    \item All objects of a class share any static member of the class.
  \end{itemize}
\end{itemize}
\end{frame}

\begin{frame}{Types of Member Functions}
  \begin{itemize}
  \item<1->
    Method \emph{overloading} -- within a class, two methods with the same name but different signatures.
  \item<2->
  Method \emph{overriding} -- same signatures across different classes (subclass and superclass).
  \end{itemize}
\end{frame}

\begin{frame}[fragile]{The \texttt{this} Pointer}
  \protect\hypertarget{pointer-this}{}
  \begin{itemize}
  \item
    Every object of a class maintains a (hidden) pointer to itself called
    \lstinline{this}.
  \item
    When an object invokes a member function, \lstinline{this} is referenced by the member function.
  \item
    Within a member function, \lstinline{this} can also be used explicitly.\\
    \lstinline{this->member;}
  \end{itemize}
\end{frame}
    

\section{Constructors}

\begin{frame}{Constructor}
  A \emph{constructor} contains code that initializes the object.

  \begin{itemize}
\item
  \emph{Default Constructor} -- Requires no parameters.
\item
  Parameterized Constructors -- Defines the object's initialization.
\item
  Constructors are executed automatically.
\end{itemize}
\end{frame}

\begin{frame}[fragile]{Constructor}
\begin{columns}[t]
  \column{0.47\textwidth}
  \texttt{AClass.hpp}
\begin{lstlisting}[basicstyle=\small\ttfamily]
class AClass {
public:
  // Default Constructor
  AClass();


  // Custom Constructor
  AClass(int size);

private:
  // A member field
  int mSize;
};
\end{lstlisting}
\column{0.5\textwidth}
\texttt{AClass.cpp}
\begin{lstlisting}[basicstyle=\small\ttfamily]
#include "AClass.hpp"

AClass::AClass()
{
  mSize = 0;
}

AClass::AClass(int size)
{
  mSize = size;
}
\end{lstlisting}
\end{columns}
\end{frame}

\begin{frame}[fragile]{Invoking a Constructor}
\begin{itemize}[<+->]
\item In Java, a constructor is invoked only through the \lstinline{new} keyword because all object variables are references.
\item In C++, a constructor is called upon variable declaration, or explicitly through \lstinline{new} with pointers, or in other situations.\\
\begin{lstlisting}[basicstyle=\ttfamily\bfseries\small]
Clock myClock; // Invokes the default constructor.
Clock *pClock = new Clock; // Invokes the default constructor
\end{lstlisting}
\item You can create and invoke custom constructors.
\begin{lstlisting}[basicstyle=\ttfamily\bfseries\small]
Clock start{12} // Invokes a constructor that accepts and int
\end{lstlisting}
\end{itemize}
\end{frame}


\section{Destructors}

\begin{frame}[fragile]{C++ Destructor}
\begin{itemize}
\item Special method whose signature is a \lstinline{~} followed by the name of the class.
\item Particularly if the class contains pointers and the constructor contains calls to new, a destructor needs to be defined.\\

\end{itemize}
% \begin{lstlisting}[basicstyle=\small\ttfamily]
% class SomeClass {
%   public:
%     SomeClass();  // Constructor
%     ~SomeClass(); // Destructor
%   private:
%     int *mpArr; // Pointer to a dynamic array.
% }
% \end{lstlisting}
% \begin{lstlisting}[basicstyle=\small\ttfamily]
% SomeClass::SomeClass() { mpArr = new int[20]; }
% SomeClass::~SomeClass() { delete [] mpArr; }
% \end{lstlisting}
\end{frame}

\begin{frame}[fragile]{C++ Destructor}
  \begin{columns}[t]
    \column{0.5\textwidth}
    \texttt{AClass.hpp}
\begin{lstlisting}[basicstyle=\small\ttfamily]
class AClass {
public:
  // Default Constructor
  AClass();

  // Destructor
  ~AClass();

private:
  // Pointer to a dynamic array
  int *mpArr;
};
\end{lstlisting}
  \column{0.47\textwidth}
  \texttt{AClass.cpp}
\begin{lstlisting}[basicstyle=\small\ttfamily]
#include "AClass.hpp"

AClass::AClass()
  : mpArr{new int[20]}
{
}

AClass::~AClass()
{
  delete[] mpArr;
}
\end{lstlisting}
  \end{columns}
\end{frame}

\section{Copying Objects}

\begin{frame}{Control Over Copy and Assignment in C++}
The \emph{copy constructor} and \emph{assignment operator} define the semantics of \lstinline{a = b}.
\begin{itemize}
  \item<2->
    The \textit{copy constructor} defines how a \emph{new} object created as a copy of an existing object (e.g., when function parameters are passed).
  \item<3->
    More low-level control is available (e.g., performing a deep copy instead of a shallow copy).
    \begin{itemize}
      \item
        Be careful to implement a deep copy if desired.
    \end{itemize}
\end{itemize}
\end{frame}

\begin{frame}{Assignment Operator}
  \begin{itemize}[<+->]
    \item
      The assignment operator (\lstinline{a = b}) where object \lstinline{a} is an existing object. Means that we are replacing the values in \lstinline{a} with a copy of the values in \lstinline{b}.

    \item
      Therefore, we must define how this should happen.

    \item
      Doing so requires \emph{overloading} the assignment operator's definition when it accepts two objects of our class.

    \item
      We will give a brief introduction to \emph{operator overloading} here but go in dept on the topic in the next lecture.
  \end{itemize}
\end{frame}

\section{Operator Overloading}

\begin{frame}[fragile]{Operator Overloading}
\begin{itemize}
\item In C++, operators like \texttt{=}, \texttt{+}, \texttt{*}, \texttt{==}, etc. can be defined, just like methods.
\item Example:

\begin{lstlisting}[basicstyle=\ttfamily\footnotesize]
class Matrix {
public:
  Matrix(const Matrix& m); // Copy constructor
  Matrix operator+(const Matrix& m); // overload + op
  Matrix operator=(const Matrix& m); // overload = op
  // ...
};
\end{lstlisting}
\begin{lstlisting}[basicstyle=\ttfamily\footnotesize]
Matrix a, b, c;
c = a + b;  // equiv to c = a.operator+(b);
\end{lstlisting}
\end{itemize}
\end{frame}

\begin{frame}[fragile]{Overloading the Assignment Operator (\texttt{=})}
  \begin{itemize}
  \item
    Function prototype:\\
    \lstinline{const className& operator=(const className&);}
  \item
    Function definition:\\
    \begin{lstlisting}[basicstyle=\small\ttfamily]
const className& operator=(const className& rightObject)
{
  // local declaration, if any
  // Avoid self-assignment
  if (this != &rightObject)
  {
    // Copy rightObject into this object
  }
  return *this; // return the object assigned
}\end{lstlisting}
  \end{itemize}
\end{frame}

\begin{frame}[fragile]{Overloading the Stream Insertion Operator
  (\texttt{{<}<})}
  \protect\hypertarget{overloading-the-stream-insertion-operator}{}
  \begin{itemize}
  \item
    Function prototype:\\
\begin{lstlisting}
friend std::ostream& operator<<(std::ostream&, const className&);
\end{lstlisting}
  \item
    Function definition:\\
\begin{lstlisting}[basicstyle=\small\ttfamily]
std::ostream& operator<<(std::ostream& out,
    const className& cObject) {
  // local declaration, if any
  // Output the members of cObject.
  // out << ...

  return out; // return the stream object.
}\end{lstlisting}
  \end{itemize}
\end{frame}

\begin{frame}{Objects as Function Parameters}
\begin{itemize}
\item<1->
  Objects can be passed as parameters to functions and returned as function values.
\item<2->
  Class objects can be passed by value or by reference.
\item<3->
  If an object is passed by value, the data members' values are copied (using the copy constructor).
  \begin{itemize}
  \item Requires additional storage space and a considerable amount of computer time
  \end{itemize}
\item<4->
  \emph{Design tip}: Use reference parameters to avoid unnecessary copies. Make the parameter \lstinline{const} if the function should not modify it.
\end{itemize}
\end{frame}



\section{Inheritance}

\begin{frame}{Inheritance vs. Composition}
Two common ways to relate two classes in a meaningful way are:
\begin{itemize}
\item
  \emph{Inheritance} -- an ``is-a'' relationship.
\item
  \emph{Composition} or \emph{aggregation} -- a ``has-a'' relationship.
\end{itemize}
\end{frame}

\begin{frame}{Inheritance}
\begin{itemize}
\item Feature that allows a class to be defined based on another class.
\begin{itemize}
\item Methods and attributes are inherited.
\end{itemize}
\item Java and C++ difference
\begin{itemize}
\item \eqmakebox[subclass][l]{Java}:  \lstinline[language=Java]{public class A extends B \{ \}}
\item \eqmakebox[subclass][l]{C++}:  \lstinline{class A: public B \{ \};} \\
(different types of inheritance)
\end{itemize}
\item Multiple inheritance possible in C++, not in Java.
\item But in Java, one may implement several interfaces.
\end{itemize}
\end{frame}

\begin{frame}{Static vs. Dynamic Binding}
\begin{itemize}
\item Let \texttt{Teacher} be a subclass of \texttt{Employee} \\
\begin{itemize}
\item Also, suppose \lstinline{promote()} is a method defined in both classes.
\end{itemize}
\item Employee variables can refer to Teachers
  \begin{tabular}{l @{\hskip 4em} l}
    \textbf{Java} & \textbf{C++}\\
    \rowcolor{codeBackground}\lstinline[language=Java]{Employee emp;} & \lstinline{Employee *emp;}\\
    \rowcolor{codeBackground}\lstinline[language=Java]{emp = new Teacher();} & \lstinline{emp = new Teacher;}\\
    \rowcolor{codeBackground}\lstinline[language=Java]{emp.promote();} & \lstinline{emp->promote();}\\
  \end{tabular}
\item Is the Employee's or the Teachers' \lstinline{promote()} method called?
\end{itemize}
\end{frame}

\begin{frame}{Static vs. Dynamic Binding}
\begin{itemize}
\item In C++, Employee's \lstinline{promote()} is called.
\begin{itemize}
\item As determined at compile time and deduced from the type of the variable (static binding).
\end{itemize}
\item In Java, Teacher's \lstinline{promote()} is called.
\begin{itemize}
\item As determined at run-time, because the actual type of the referred object is checked then (dynamic binding)
\end{itemize}
\end{itemize}
\end{frame}

\begin{frame}[fragile]{Static vs. Dynamic Binding}
\begin{itemize}[<+->]
\item
  A function can be declared \textit{virtual} to use dynamic binding instead of static.\\
  \begin{lstlisting}
class Animal {
public:
  virtual void print() const;
  // ...
};\end{lstlisting}
\item
  With inheritance, it is a good idea to make the destructors virtual to ensure that the child object's memory is freed.
\end{itemize}
\end{frame}

\begin{frame}{Design Principles}
\begin{itemize}
\item
  \emph{Encapsulation}: combines data and operations on data in a single unit.
\item
  \emph{Inheritance}: creates new objects (classes) from existing objects (classes).
\item
  \emph{Polymorphism}: the ability to use the same expression to denote different operations.
\end{itemize}
\end{frame}

\begin{frame}{Static vs. Dynamic Polymorphism}
\begin{itemize}
\item
  Inheritance provides the ability to change types at runtime (though subclasses).
\item
  Templates provide the ability to change types at compile time.
  \begin{itemize}
  \item
    We will learn more about templates in the next set of slides (Chapter 13).
  \end{itemize}
\end{itemize}
\end{frame}


\begin{frame}{There is more.}
\begin{itemize}
\item
  We have quickly skimmed the material in Chapters 10 -- 11 in these slides.
\item
  Reference the textbook for details and examples.
\end{itemize}
\end{frame}

\end{document}
